\paragraph{Computing orders of vanishing on the flag.}
Work over the deep coefficient field \(F\). Let \(v\) be a
discrete valuation on \(F(X_Z)\) centred at the boundary point
\(y\), including a divisorial valuation on a modification whose
centre on \(X_Z\) has arbitrary codimension.

\begin{lemma}[A common valuation ring for boundary evaluations]\label{lem:common-boundary-trait}
The marked generic flag admits a split trait with closed point
over \(y\). For its complete DVR \(V\), normalized order \(v_V\)
and local map \(g:\mathcal O_{X_Z,y}\to V\), one has a common proportionality constant
\[
 \exists\,c>0\quad \forall\,a\in\mathcal O_{X_Z,y}\setminus\{0\},\qquad v(a)=c\,v_V(g(a)).
\]

\end{lemma}
\begin{proof}
Complete the valuation ring dominating \(\mathcal O_{X_Z,y}\).
The generic flag remains in the open marked locus. Pull back the
semiabelian family and choose an integral eigenframe adapted to
the flag by triangular change of basis, together with a generating
theta frame. Extend to split the reduction torus. Polarized
uniformization gives a period lattice \(X\), an integral quadratic
form \(Q\) and a balanced pairing with \(\beta(x,x)=2Q(x)\);
the adjoint on \(\mathcal O_E\) is CM conjugation. If \(e'>0\)
is the ramification index of the extension and the initial
comparison is \(v=c_0v_{\rm old}\), then
\[
             v_V|_{\rm old}=e'v_{\rm old},\qquad c=c_0/e' .
\]
The map preserves the generic flag and sends the closed point
to \(y\). The structure map over \(F\) embeds \(F\) in the
residue field, which therefore has characteristic zero.
\label{cor:deep-boundary-orders-1}
\end{proof}

\begin{lemma}[Normalization of evaluation in a boundary frame]\label{lem:deep-trait-frame-order}
At the selected valuation centred on the boundary, use the
split trait and comparison factor \(c>0\) of Lemma~\ref{lem:common-boundary-trait}. Define \(\phi=c\inf_h h(\beta)\), where \(\beta\) is the
trait's monodromy form. The infimum runs over real linear functionals
on symmetric balanced forms that take integral values on integral
forms and are strictly positive on every nonzero positive semidefinite
form. For a nonzero germ \(a\) of any weight line of central theta exponent zero, write \(a_V\) for its value in the integral frame over the valuation ring. In the local line frame,
\[
 \phi=c\inf_h h(\beta)>0,\qquad v(a)=c\,v_V(a_V).
\]

\end{lemma}

\begin{proof}
\proofstep{The monodromy parameter on the common trait.}
For the split trait define \(\phi_V\) to be the infimum of
\(h(\beta)\) for real linear functionals \(h\) on the space
of symmetric balanced forms which are integral on integral
forms and strictly positive on every nonzero positive
semidefinite form. The toric rank is positive. Otherwise
the abelian special fibre would place the image of the
closed point in the model interior. The polarization gives
positive definite \(Q\) on its nonzero period lattice.
Thus \(\beta\) is nonzero and integral. Strict positivity and integrality give
\[
 h(\beta)\in\mathbb Z_{>0}\quad\text{for every admissible }h,\qquad \phi_V\ge1.
\]
The admissible
set is nonempty by the spanning construction below. Set
\(\phi=c\phi_V>0\).

In each trivialization of a weight line at \(y\), the
coefficient of an open-locus rule, evaluated at the generic
point of this trait, is a unit of \(V\) times the integral
frame evaluation. Indeed the adapted eigenbases are related
by a triangular flag change whose diagonal factors are units
because both generate the successive integral quotients.
If \(u_i\) are the diagonal units and \(a_i\) the integral exponents of the weight character, the multiplier is a unit because
\[
 v\!\left(\prod_i u_i^{a_i}\right)=\sum_i a_i v(u_i)=0.
\]
The central
theta exponent is zero, so the result is a unit also for
negative weight entries. It has zero trait order, and the
only change of valuation is the common factor \(c\). It applies equally to the
coefficient of a section extending the rule. In particular
a nonzero value in the frame gives a nonzero germ, with order \(c\)
times the frame order. \label{cor:deep-boundary-orders-2}

\end{proof}

\paragraph{A single section with large order.}
The choice of the section begins on the geometric generic fibre. We
need a product of nonzero theta values there, while retaining the
residues of its labels modulo five for the boundary calculation.

Take \(R=5\). The deep principal level contains \(R^4\), since
\(\lfloor3L/\log3\rfloor\ge4\). Write \(P_x\) for the torsion
point with label \(x\) in the full \(R^4\)-marking. The null
value of order \(R^2\) at \(P_x\) is the theta section at
\([R^2]P_x\), expressed in the fibre trivialization given by the
theta structure. The trivialization is an isomorphism. For labels
\(x+R\widetilde u\), with \(u\) varying modulo \(R\), we have
\[
   [R^2]P_{x+R\widetilde u}=[R^2]P_x+P_{R^3\widetilde u},
   \qquad \{P_{R^3\widetilde u}:u\in\Lambda/R\Lambda\}=A[R],
\]
Here \(\Lambda\) is the marked period lattice.

The geometric
generic field embeds in \(\mathbb C\) since \(Z\) is of finite
type over a number field. In this embedding
Lemma~\ref{lem:eta-torsion-translate} gives a shift with a nonzero
value for every label. The full marking defines each shifted
evaluation over the geometric generic field, so nonvanishing descends.
For a generating tuple \((x_b)_{b=1}^N\) choose the shifts
independently. If \(A\) denotes the normalized theta value, then
\[
 f=\prod_{b=1}^N A(x_b+R\widetilde u_b)\ne0,\qquad
 x_b+R\widetilde u_b\equiv x_b\pmod{R\Lambda}\quad(1\le b\le N).
\]
This single section is nonzero on the component. The congruence
retains the generating property in every split uniformization.
\label{cor:deep-boundary-orders-3}

\begin{lemma}[One auxiliary section on each component]
\label{lem:deep-boundary-uniform-order}
On a component \(Z\) of the marked compact flag, there is a
single nonzero auxiliary section \(\theta\) such that for each allowed centered boundary valuation \(v\), with normalized boundary order \(\phi_v\),
\[
 v(\theta)\ge\frac{8}{625}\phi_v.
\]

The section is fixed before the center or comparison trait is
chosen.
\end{lemma}

\begin{proof}
\proofstep{The auxiliary lower order from the half-level product.}
Here is the lower order estimate. Put \(m=5^2\). In a
Heisenberg lift at level \(m^2\), let \(y_j\in X\) be
the period coordinates of the shifted labels. They generate
\(X\) modulo \(m^2X\) as an \(\mathcal O_E\)-module.
The integral odd-level expansion of the normalized
theta-null at the \(j\)-th label has order at least
\[
             m^{-2}\min_{z\in X}Q(y_j+mz).
\]
This is the valuation bound term by term in the theta
expansion on that coset: the normalized quadratic term
has the indicated order and the integral coefficient has
nonnegative order. The valuation can be extended to the
algebraic closure for this computation. The product is
the value of \(f\) over the original fraction field.

Choose minimizing representatives \(b_j=y_j+m z_j\).
The quotient of \(X\) by their \(\mathcal O_E\)-span
is finitely generated over \(\mathbb Z\) and is equal
to its multiple by \(m\), so it has rank zero. Take the
marked half-trace element \(w_0\in\mathcal O_E\) with
\(w_0+\bar w_0=1\). Then
\[
 h(\gamma)=\sum_j\gamma(w_0 b_j,b_j)
          =\tfrac12\sum_j\gamma(b_j,b_j)
\]
is integral on integral balanced forms. For a semidefinite
\(\gamma\), equality to zero puts every \(b_j\) in its
radical. Balance and finite index then put all of \(X\)
in its radical. Hence \(h\) is admissible and
\(\sum_j Q(b_j)=h(\beta)\ge\phi_V\). The product has
order at least \(\phi_V/625\).

\proofstep{Transporting the fixed section to every centre.}
Use the key rule \(K\) and the component flat rules
\(s_\tau\) of the preceding flat-line construction. On
\(Z\) the required rule is
\[
        f^8 K^N \prod_\tau s_\tau^{4N}.
\]
Extension by the component idempotent specifies its support.
The key cancels the central exponent, and the flat power
restores weight \(8N\omega_{\rm H}\). All factors after
\(f^8\) are invertible there. The key on a generating
integral theta and Lie frame is integral, and the flat rules have order zero here. More explicitly,
the rank-one and de Rham canonical extension comparisons
give a horizontal generator taking a unit value in the
integral frame on a trait over \(F\). The constant-length
rule of the same weight differs by a scalar in \(F^\times\)
on this component. For such a scalar \(a\) the inclusion \(F\subset V\) gives
\[
 a,a^{-1}\in V,\qquad v(a)=0.
\]

Thus every such nonzero value has order at least
\(8\phi_V/625\). The rule defines a rational section,
regular on the open. At every codimension-one boundary
point the calculation with its local DVR gives nonnegative
order. Normality extends the rule to a section \(\theta\)
on \(X_Z\). It remains nonzero by its value on the generic fibre.
The same calculation at each centred valuation now gives
\[
 v(\theta)\ge 8\phi/625.
\]

\label{cor:deep-boundary-orders-4}

\end{proof}

\paragraph{A cusp section with small order at a given centre.}
We give the weight and coordinate details because this
section varies with the valuation. For the selected quotient weight \(u_1\), write \(p,q\ge1\) for the two block sizes, \(p+q=r\). Let \(q_\tau\) be the size at an embedding of \(E\), equal to \(p\) or \(q\) on the two embeddings of each conjugate pair. Define
\begin{equation}\label{eq:deep-quotient-column-count}
 c=\sum_{\tau}\sum_{j=1}^{q_\tau}2j
   =f[p(p+1)+q(q+1)],\qquad
 n=8\big(\lfloor c/8\rfloor+fr+1\big).
\end{equation}
There are \(fr\) nonempty row labels. Hence \(c\le n\le2fr^2+8fr+8\), and the \(j\) columns for every \(j\)-row minor may be selected disjointly across all labels. This is the column number for the quotient-weight differential.

At \(\tau\) in the chosen CM type, write the decreasing coordinates of \(u_1\) as \(u_{\tau,1},\ldots,u_{\tau,r}\). Let \(h_0\) be the undivided inter-block gap and \(d_{\rm fl}=f\binom r2\) the flag dimension. The gap reserve in the weight selection gives
\[
 4Tn\le80Tfr^2\le T r f^2m^2 k d_{\rm fl}\le h_0 .
\]

Euclidean division at length \(T\) gives \(4T(u_{\tau,p}-u_{\tau,p+1}-3r)\ge h_0\), so \(u_{\tau,p}-u_{\tau,p+1}\ge n+3r\). Use the parallel shift \(w_\tau=n/2+1-u_{\tau,p}\) and central exponent \(n/2\). Their polynomial exponents, in the original and conjugate block, are
\begin{align}
 a_{\tau,j}&=u_{\tau,j}+1-u_{\tau,p} &&(1\le j\le p),\nonumber\\
 a_{\bar\tau,j}&=u_{\tau,p}-u_{\tau,r+1-j}-n-1
                                    &&(1\le j\le q).\label{eq:deep-quotient-exponents}
\end{align}
Thus the two arrays are decreasing and nonnegative. The last exponent on the chosen side is one. Each exponent is at most \(4b+1\) because \(|u_{\tau,j}|\le2b\), so the total degree \(D_a\) is at most \(fr(4b+1)\).

Put
\[
 H=2fm,\qquad R_E=|\operatorname{disc}E|^{2/[E:\mathbb Q]}
                    /(2[E:\mathbb Q]),\qquad M=g_0g_1,
\]
where \(m\) here denotes the rank parameter in the original family. The
original half-level generator satisfies
\(\alpha\bar\alpha=M\), the deep level is divisible by
\(M\) and its marked half ideal is contained in
\((\bar\alpha)\). The two half ideals are coprime, so
\(\alpha,\bar\alpha\) generate the unit ideal. These are
the level conditions for translating a jet at level \(M\)
inside this cover.

The numerical hypothesis used at \(u_1\) is
\[
 4D_a+4\le g_0,\quad 2D_a\sqrt H<g_0,\quad
 \frac{R_E H^3(rH/4+2D_a^2)}{g_0}\le\frac1{2500}.
\]
Indeed the quotient-weight and size inequalities give
\(D_a\le fr(4b+1)\), \(f,m,r\le t\), \(t\ge11\),
\(t^7\le b\), \(R_E\le b\), \(b^8\le g_0\).
They imply \(D_a\le b^2\), \(H^3\le b\),
\(rH\le b^2\). The numerator is at most \(3b^6\)
and \(b\ge100\), which proves the last inequality.
The first two follow from \(4b^2+4<b^8\) and
\(2b^2\sqrt H<b^8\).
\label{cor:deep-boundary-orders-5}

\begin{lemma}[A cusp section detecting a prescribed center]
\label{lem:deep-boundary-detecting-section}
For each allowed centered boundary valuation there is a cusp
section at the selected quotient weight \(u_1\) whose value \(a\) in the chosen frame satisfies, in the original normalization,
\[
 a\ne0,\qquad v(a)\le\phi/2500.
\]
 This section may depend on the center. The
constant and size thresholds are independent of the center,
component and ramification.
\end{lemma}

\begin{proof}
\proofstep{The detecting jet at the chosen centre.}
Use a truncated integral frame chart, after a further
finite extension if necessary. Denote its ramification
index by \(e_1\), its boundary parameter by
\(\Phi=e_1\phi_V\), and its toric rank by \(s\le H\).
For clarity, the local theta-jet estimate used here is the
following one. For a truncated chart with split polarized
uniformization and torsion identification, residue
characteristic zero, an antitone exponent array and an even column count divisible
by eight and at least twice the sum of row ranks, the infimum of
normalized orders of its marked harmonic theta translates
is at most
\[
 \nu_{\rm loc}(P)+D_a^2 R_E s^3\Phi/M .
\]
Here \(P\) is a harmonic polynomial of this weight and
\(\nu_{\rm loc}(P)\) is the infimum of its contracted jet
orders on integral germs. The estimate requires the half
generator just specified, \(4D_a+4\le M\) and
\(2D_a\sqrt s<M\).

To see the loss, take the period basis \(b_j\) given by
reduction of the balanced Gram form. It satisfies
\[
 Q(b_j)\le H_0=R_Es^3\Phi,\qquad
 |Q(b_j+z)-Q(b_j)-Q(z)|\le 2\sqrt{s}\,Q(z)
 \quad(z\in X).
\]
On the product of \(n\) copies, place one basis vector
in one coordinate at a time. The form on the product is
the sum of the forms on the coordinates, so the same
two bounds hold without an extra factor of \(n\).
To lift an integral jet \(J\) through degree \(D_a\), take
the small characters \(\chi=\sum_j n_j b_j\) on this product
with \(n_j\ge0\) and \(\sum_j n_j\le D_a\).
The triangle inequality for \(\sqrt Q\) gives
\(Q(\chi)\le D_a^2 H_0\).

For a nonzero period \(z\),
\[
 Q(\chi+Mz)-Q(\chi)
     \ge M(M-2D_a\sqrt{s})Q(z)>0 .
\]
Thus \(\chi\) is the unique minimum in its level coset.
This shows why reduction of the corresponding integral
theta expansion singles out the \(\chi\)-term. At
\[
 M\ge4D_a+4
\]
the integral abelian jet theorem permits
arbitrary base jets in its coefficient line, in a local
generator. The remaining torus jets are spanned over those
base jets by the small characters: the differences of the
basis characters from their values at the identity generate
the torus augmentation directions, and products of at most
\(D_a\) such differences suffice in the truncated ring.
Frame factors here are units. Reduction therefore spans
the whole jet module on the special fibre. This finite
truncated module over the local ring is spanned by the lifted
jets by Nakayama. Write \(J\) as an integral combination
of them.

In their generic expansions the valuation loss
per character is at most \(Q(\chi)/M\). Over a valued
algebraic closure choose \(\lambda\) of order
\[
 \lfloor D_a^2 H_0\rfloor/M
\]
(a root of a uniformizer
power suffices). Since \(Q(\chi)\) is integral, this
one scalar clears all the losses. The resulting global
expansion has nonnegative normalized coefficient orders
and jet \(\lambda J\).
The coprime half generators
and the unit \(M\) in the residue field allow the
torsion-translate spanning formula for this expansion.
Its coefficients \(c_j\) have nonnegative valuation.
Uniqueness of expansion identifies its germ with the
sum of the translated germs in the same frame. Thus,
for harmonic contraction \(P(\nabla)\), which only sees
the jet through degree \(D_a\), one has
\[
 \lambda P(\nabla)J=\sum_j c_j P(\nabla)J_j.
\]
A nonzero left side gives a nonzero summand whose
unweighted order is at most
\(v(P(\nabla)J)+D_a^2H_0/M\).

In passing from the
translate to its marked rule one removes a multiplier
of nonnegative order. The canonical connection and the
level factorization identify the harmonic contraction on
both sides. Taking the infimum over integral germs proves
the estimate. Here orders of evaluated rules are taken in
the chart DVR even though the translation identity can be
formed over its algebraic closure.
\label{cor:deep-boundary-orders-6}

\proofstep{The compatible highest-weight polynomial.}
Choose for \(P\) the product of successive highest-weight
minors with distinct assigned columns. Dominance gives
nonnegative difference exponents and the disjoint columns
for conjugate rows give harmonicity. On the saturated
chart directions the dual monomial jet gives a nonzero
product of factorials as contraction. These integers
are units in residue characteristic zero. Consequently
\[
 \nu_{\rm loc}(P)\le0.
\]
The bound for the infimum
is now at most \[
 B=D_a^2R_Es^3\Phi/M\le\Phi/2500.
\]
A zero evaluation has infinite order. Since the chart's order has normalization factor one, every nonzero value \(a\) satisfies
\[
 v(a)>B\quad\Longrightarrow\quad v(a)\ge\lfloor B\rfloor+1.
\]
If all nonzero values satisfied the left inequality, their infimum would satisfy the same integer lower bound. Thus there is
a rule with a nonzero value of order at most \(B\).

\proofstep{Returning its integer order to the original trait.}
For the same theta rule \(f\) on the extended chart and original integral frame the two orders satisfy
\[
 v_{\rm chart}(f)=e_1 v_{\rm frame}(f).
\]
Restriction to the component gives a cusp rule
with the same value: the dominant exponents, last positive
entry and column range are those of the harmonic cusp
construction. Its weight differs from \(u_1\) by a parallel
shift and central exponent \(n/2\). Multiply by the
inverse parallel twist and the inverse key twist for this
central degree. The exponent needed for the weight return is integral, since
\[
 4\mid n/2\quad\Longrightarrow\quad n/8\in\mathbb Z.
\]
Their evaluations are units here,
including inverse powers since the component idempotent
reads one.

The cusp extension to the compact flag therefore
provides a section at \(u_1\) with a nonzero value in the frame of order at most \(\phi_V/2500\). Multiplication by the
original \(c\) bounds the order of its chart germ by
\(\phi/2500\). Every choice after fixing the auxiliary
section may depend on the centre. The constants and
the family size thresholds used for these two inequalities
are independent of that choice, the component, and the
ramification of either trait.
\label{cor:deep-boundary-orders-7}
\end{proof}
